\chapter{Initial Implementation Analysis}
\label{chap:initial-implementation}

This chapter analyzes the implementation that was inherited at the beginning of the project.
The objective is not yet to discuss the optimized variants introduced later, but to explain
what the starting system already did, how its components interacted, and which practical
limitations motivated the rest of the work. This distinction matters because the current
repository now contains both inherited components and new specialized contracts. Unless
explicitly stated otherwise, this chapter refers to the initial execution path centered on
the original optimistic account, the original dispute account, the Rust/WebAssembly
precontract pipeline, the Next.js backend, the local SQLite database, and the Alto/EntryPoint
transaction path.

\section{Repository structure and technology stack}

The initial codebase was not a single smart contract repository. It was a full proof of
concept combining off-chain computation, on-chain adjudication, a local web application, and
transaction-sponsorship infrastructure. The main components are summarized in
Table~\ref{tab:initial-repo-components}.

\begin{table}[htbp]
  \centering
  \begin{tabularx}{\linewidth}{>{\raggedright\arraybackslash}p{0.28\linewidth}>{\raggedright\arraybackslash}X}
    \toprule
    Component & Role in the initial implementation \\
    \midrule
    \path{src/wasm} &
    Rust crate used for cryptographic preprocessing, circuit construction, ciphertext
    generation, commitments, and native command line tools. It also produces the
    WebAssembly bundle used by the application. \\
    \path{src/hardhat} &
    Hardhat workspace containing Solidity contracts, deployment scripts, generated
    artifacts, and tests. The initial path uses an optimistic account and a dispute account,
    plus libraries for accumulators, commitments, AES, SHA-256, and circuit evaluation. \\
    \path{src/app} &
    Next.js application containing the frontend, API routes, SQLite database access,
    blockchain helper code, and local file storage endpoints. \\
    \path{bundler-alto} &
    Local Alto bundler used to submit UserOperations to the local chain through an
    EntryPoint-compatible account-abstraction path. \\
    \path{desktop} &
    Optional Electron shell around the Next.js application. It is useful for local
    demonstrations and for executing native preprocessing outside the browser. \\
    Root scripts &
    Shell scripts such as \path{install.sh}, \path{deploy-all.sh}, and
    \path{run-alto.sh} orchestrate installation, contract deployment, and bundler startup. \\
    \bottomrule
  \end{tabularx}
  \caption{Main components of the inherited implementation.}
  \label{tab:initial-repo-components}
\end{table}

The local development stack is based on a Hardhat chain, Solidity contracts compiled through
Hardhat, a Next.js server running on \path{localhost:3000}, a bundler RPC endpoint running on
\path{localhost:4337/rpc}, and a SQLite database stored under \path{src/app/db}. This setup
is appropriate for a research prototype because it makes the complete execution path visible
and reproducible. It is not, however, equivalent to a production deployment: identities are
local test accounts, the backend is centralized, and the bundler is a local development
service.

The contract compiler configuration also reflects this research-prototype status. The
Hardhat configuration enables optimization and compilation through the Solidity intermediate
representation, and local networks are configured with unlimited contract size. This was
necessary because the original dispute contracts and deployer libraries were large. It also
means that local deployment success alone cannot be interpreted as mainnet deployability.

\section{Rust and WebAssembly off-chain pipeline}

The most important off-chain component is the Rust crate in \path{src/wasm}. It implements
the preprocessing required before an exchange can be proposed. The native binary
\path{precontract_cli} takes an input file and an optional 16-byte AES key. If no key is
given, it samples a random one. It then computes the precontract material used by the rest of
the system.

The output of this command line tool contains, in particular:

\begin{itemize}
  \item the description output encoded as hexadecimal data;
  \item the ciphertext commitment \(h_{ct}\);
  \item the circuit commitment \(h_{circuit}\);
  \item the commitment value \(c\) and opening value \(o\);
  \item the number of ciphertext blocks and the number of circuit gates;
  \item paths to the generated ciphertext and circuit files;
  \item the AES key, returned so that the vendor can later reveal it on-chain.
\end{itemize}

The same crate contains the circuit construction and evaluation code. The implementation uses
a compact V2 circuit representation based on fixed 64-byte gates. The relevant operations
include AES-CTR decryption gates, SHA-256 compression gates, constants, XOR operations, and a
final comparison gate. This representation is central to the dispute mechanism because it
lets the smart contract evaluate a single local gate once the corresponding gate bytes and
input values are authenticated.

The Rust code is exposed in two ways. The WebAssembly bundle is used by the web application
for interactive use, while native command line tools are used by the backend and by
performance experiments. This dual execution path was already important in the inherited
system. The browser/WebAssembly path is convenient for demonstrations and moderate files,
but native execution is more appropriate for large files because it avoids browser memory
and UI constraints.

The inherited proof-generation API was only partial. The route
\path{src/app/api/proofs/compute/route.ts} calls the native \path{compute_proofs_cli}
binary, but the exposed path supports only one dispute state in the current application
interface: the state corresponding to \texttt{WaitVendorDataRight}. This was sufficient for
some local tests and demonstrations, but it did not provide a complete API for all possible
Step 8 cases. This limitation is one reason why later measurements use dedicated test
scripts and native benchmark binaries rather than relying exclusively on the web API.

\section{Frontend, backend, database, and file storage}

The user-facing application is implemented with Next.js and a small set of role-oriented
components. The user side contains views for creating contracts, accepting precontracts, and
following ongoing contracts. The sponsor side contains views for contracts waiting for
sponsorship and for dispute sponsorship. The application also has blockchain helper modules
for optimistic contracts, dispute contracts, UserOperations, deployment configuration, and
contract artifacts.

The backend stores protocol metadata in SQLite. The main table, \texttt{contracts}, stores
the precontract data needed by the application: item description, commitment opening,
buyer and vendor addresses, price, number of blocks, number of gates, commitment,
completion and dispute tips, protocol version, timeout delay, algorithm suite, acceptance
status, sponsor address, deployed optimistic contract address, and session-key information.
A second table, \texttt{disputes}, links a contract identifier to buyer-side and vendor-side
dispute sponsors and to the deployed dispute contract address.

The encrypted file itself is stored off-chain under the application upload directory. During
precontract creation, the backend can receive a multipart upload, call the native
\path{precontract_cli}, copy the generated ciphertext to \path{src/app/uploads}, and insert
the corresponding metadata into SQLite. Alternatively, in Electron mode, the frontend can
send precomputed values produced locally.

This design is practical for a proof of concept because it makes the whole system easy to
run on one machine. Cryptographically, the backend is not trusted to decide the fairness of
the exchange: the smart contract only relies on commitments and proofs. Practically, however,
the backend remains an availability and orchestration dependency. If the ciphertext or proof
artifacts are not available when needed, the on-chain protocol cannot reconstruct them by
itself. This is not a contradiction of the protocol, but it is an important implementation
assumption.

The initial application also uses predefined local accounts. Public/private key pairs are
hard-coded for the local Hardhat chain, and the UI behaves like a laboratory multi-user
interface rather than a wallet-authenticated application. This makes testing much faster, but
it must not be confused with a real wallet integration.

\section{Optimistic smart contract architecture}

The initial optimistic phase is implemented by an account-style contract. The central
contract in the inherited execution path is \path{OptimisticSOXAccount.sol}. It combines
three responsibilities that are conceptually distinct:

\begin{itemize}
  \item it stores the exchange parameters and optimistic state;
  \item it escrows payments, tips, and sponsor funds;
  \item it also behaves as an account-abstraction compatible account for sponsored calls.
\end{itemize}

The optimistic state machine contains the states
\[
  \texttt{WaitPayment},\quad
  \texttt{WaitKey},\quad
  \texttt{WaitSB},\quad
  \texttt{WaitSV},\quad
  \texttt{InDispute},\quad
  \texttt{End}.
\]
The normal execution begins in \texttt{WaitPayment}. The buyer calls
\texttt{sendPayment}, after which the contract waits for the vendor to reveal the AES key
with \texttt{sendKey}. Once the key is available, the buyer can either complete the exchange
with \texttt{completeTransaction} or move toward a dispute by registering the buyer-side
dispute sponsorship. The vendor-side dispute sponsorship then creates or activates the
dispute contract and moves the optimistic contract into \texttt{InDispute}.

The constructor of the optimistic account receives the EntryPoint address, vendor and buyer
addresses, price, tips, timeout increment, commitment, number of ciphertext blocks, number of
gates, and vendor signer. The inherited design deploys a new optimistic account for each
exchange. This is simple and modular: each exchange has its own state, its own balance, and
its own account-abstraction context. The cost is that deployment becomes part of the marginal
cost of every exchange.

The contract also implements timeout logic. If the expected party does not act before the
deadline, \texttt{cancelTransaction} redistributes funds according to the current state. This
is important for liveness: an honest party should not be blocked forever by a counterparty
who stops responding.

\section{Dispute smart contract architecture}

The dispute phase is implemented by a separate dispute account. In the inherited path, the
dispute account is deployed when the optimistic contract reaches the dispute stage. The main
contract is \path{DisputeSOXAccount.sol}, supported by several verifier and evaluator
libraries.

The dispute contract keeps the state of the binary search. Its states cover buyer
challenges, vendor opinions, vendor data submissions, completion, cancellation, and final
termination. In the Solidity implementation these states are named
\texttt{ChallengeBuyer}, \texttt{WaitVendorOpinion}, \texttt{WaitVendorData},
\texttt{WaitVendorDataLeft}, \texttt{WaitVendorDataRight}, \texttt{Complete},
\texttt{Cancel}, and \texttt{End}. The buyer submits a response with
\texttt{respondChallenge}. The vendor answers with \texttt{giveOpinion}, which updates the
search interval. When the search reaches a local position or a boundary case, the vendor
submits authenticated local data through one of the \texttt{submitCommitment} calls.

The generic Step 8 verification is expensive because it does several things at once. It
opens the commitment to recover \(h_{circuit}\) and \(h_{ct}\), checks Merkle-style proofs
for the relevant circuit gate and values, evaluates the local V2 gate, and checks the
consistency of the resulting accumulator value. The contract therefore depends on helper
contracts and libraries for accumulator verification, commitment opening, V2 circuit
evaluation, AES-CTR evaluation, SHA-256 evaluation, and simple operations.

This architecture follows the intended optimistic principle: the whole file and the whole
circuit are never evaluated on-chain in the honest case. Nevertheless, once a dispute
starts, the deployment and local verification costs are substantial. The initial
implementation also treats the circuit generically, so the contract must be given the
64-byte gate encoding and the proof that this gate belongs to \(h_{circuit}\). This generic
design is flexible, but it is not optimal when the description is known to be
\(\mathsf{SHA256}\) and the circuit structure is determined by the input length.

\section{ERC-4337, Alto bundler, and EntryPoint integration}

The inherited implementation does not only implement protocol sponsorship at the economic
level. It also integrates a transaction-sponsorship path inspired by ERC-4337. The local
stack deploys or configures an EntryPoint-compatible contract and starts an Alto bundler. The
frontend builds UserOperations, sends them to the bundler RPC endpoint, and waits for the
corresponding UserOperation receipt.

Both the optimistic account and the dispute account expose account-style functionality:
signature validation, nonces, execution helpers, and EntryPoint deposit management. This
allows actions such as payment, key release, challenge response, and dispute evidence
submission to be executed through the account-abstraction path. In the local prototype, this
is useful because it demonstrates how a participant can interact with the protocol without
manually sending every transaction as a regular externally owned account transaction.

This integration also introduces additional assumptions. The bundler must be running, the
EntryPoint address must match the deployed configuration, and the UserOperation must be
accepted and included. From the point of view of protocol analysis, the bundler is not part
of the original fairness proof. It is an implementation layer that affects how gas is paid
and how liveness is achieved in practice. This is why later chapters separate protocol
variants from transaction-sponsorship infrastructure.

\section{Initial end-to-end execution flow}

The initial end-to-end flow validated during the first phase of the project is shown in
Table~\ref{tab:initial-e2e-flow}. It starts before any smart contract call, with the vendor
creating the precontract material, and ends either in optimistic completion or in a dispute
contract.

\begin{table}[htbp]
  \centering
  \begin{tabularx}{\linewidth}{>{\raggedright\arraybackslash}p{0.13\linewidth}>{\raggedright\arraybackslash}p{0.25\linewidth}>{\raggedright\arraybackslash}X}
    \toprule
    Order & Location & Action \\
    \midrule
    1 &
    Rust / backend &
    The vendor uploads or selects a file. The preprocessing pipeline encrypts it, builds the
    circuit, computes \(h_{ct}\), \(h_{circuit}\), and the commitment, and stores the
    generated artifacts. \\
    2 &
    SQLite / uploads &
    The backend records the precontract metadata and stores the encrypted file off-chain. The
    buyer can see pending precontracts addressed to its public key. \\
    3 &
    Frontend &
    The buyer accepts the precontract. In the initial application this is a database-level
    acceptance action before the on-chain optimistic account is deployed or sponsored. \\
    4 &
    Solidity / Hardhat &
    The optimistic account is deployed with the precontract parameters and sponsor-related
    funds. \\
    5 &
    Optimistic account &
    The buyer pays into the contract, then the vendor sends the AES key. If the buyer accepts
    the decrypted item, the exchange completes. \\
    6 &
    Dispute path &
    If the buyer complains, buyer-side and vendor-side dispute sponsorship are registered and
    a dispute account is deployed. The dispute then proceeds through challenge rounds and
    Step 8 verification. \\
    \bottomrule
  \end{tabularx}
  \caption{Initial end-to-end execution flow observed in the inherited implementation.}
  \label{tab:initial-e2e-flow}
\end{table}

This flow was essential to reconstruct before changing the protocol variants. Several gas
costs depend on decisions made outside the Solidity functions themselves. For example, a
measurement can include contract deployment, EntryPoint deposit funding, database-level
precontract acceptance, dispute deployment, or only a single function call. The initial
analysis therefore established the discipline used later in the report: every measurement
must say which part of this flow it covers.

\section{Limitations of the initial implementation}

The inherited implementation was already a complete and useful proof of concept, but it had
several limitations for the goals of this project.

\paragraph{Deployment cost was paid per exchange.}
The initial architecture deploys an optimistic account for each exchange and a dispute
account when a dispute starts. This makes the state isolation simple, but it makes deployment
cost part of the marginal cost of the protocol. For gas optimization, this is a major issue:
a large fixed deployment cost can dominate the cost of an honest exchange.

\paragraph{The optimistic and dispute accounts were too general.}
The initial contracts were designed to support the full sponsored protocol. This is useful
for correctness and demonstration, but it means that exchanges pay for code paths that are
not needed in simpler variants. For example, the normal sponsored case, no \(S\)-deposit,
\(S=B\), \(S=V\), self-sponsored disputes, and hardcoded \(\mathsf{SHA256}\) do not all need
the same storage fields and runtime branches. This observation motivated the later move from
monolithic variant support to specialized contracts.

\paragraph{The generic dispute verifier pays for circuit flexibility.}
In the generic dispute path, Step 8 authenticates a supplied 64-byte gate against
\(h_{circuit}\). This is necessary for arbitrary circuits, but it is redundant when the
circuit is fully determined by public parameters. The \(\mathsf{desc}=\mathsf{SHA256}\)
case is the main example: once the input length is fixed, the sequence of SHA-256 circuit
gates is determined. This created the opportunity for the hardcoded SHA-256 dispute mode
implemented later.

\paragraph{The backend remains an orchestration and availability dependency.}
The encrypted file and auxiliary artifacts are stored off-chain. The commitments prevent the
backend from silently changing the data without detection, but they do not make unavailable
data appear. A robust deployment would need a clearer availability layer, for example
replicated storage or explicit user-side artifact management.

\paragraph{The web proof-generation path was incomplete.}
The backend API for proof computation exposed only a limited part of the dispute proof
generation flow. Complete experimental coverage therefore required Hardhat tests and native
Rust benchmark tools rather than relying only on the interactive web interface.

\paragraph{The identity model was a laboratory model.}
The application uses predefined Hardhat accounts and a local mapping from public addresses to
private keys. This is appropriate for reproducibility, but it does not provide real wallet
authentication or production-grade key management.

\paragraph{Large files required the native path.}
The Rust/WebAssembly architecture made preprocessing portable, but the browser path was not
the right environment for very large files. For large-file measurements, especially the 1 GiB
experiments discussed later, the native Rust command line tools were the reliable execution
path.

These limitations do not invalidate the initial implementation. On the contrary, they define
the engineering problem addressed by this project. The inherited system already made the
protocol executable end-to-end; the remaining question was how to restructure it so that the
important protocol variants could be measured at the correct scope and with lower gas cost.
